% ==================================================================
% id:       func-018
% topic:    פונקציות
% subtopic: פונקציה הפוכה
% solution: full
% ==================================================================
\textbf{בדקו את ההפיכות של הפונקציות הבאות. מצאו פונקציות הפוכות בתחומים בהם קיימת פונקציה כזו.
כמו כן ציינו באילו תחומים הפונקציות מונוטוניות עולות או יורדות:}
\begin{enumerate}\bMs{2}
\item
$f(x)=1-x^3$
\begin{solbox}
הפונקציה
$f(x)=1-x^3$
חד-חד ערכית ועל
$\mathbb{R}$,
ולכן הפיכה.
בדיקה ישירה מראה כי
$$f(x)=y \iff 1-x^3=y \iff x^3=1-y \iff x=\sqrt[3]{1-y}$$
לכן
ההופכית היא
\[
f^{-1}(x)=\sqrt[3]{1-x}.
\]
הפונקציה מונוטונית יורדת ממש בכל $\mathbb{R}$.
ניתן לראות זאת ישירות כי
$$.x_1<x_2\Rightarrow x_1^3<x_2^3\Rightarrow 1-f(x_1)=x_1^3> 1-x_2^3=f(x_2)$$
\end{solbox}
\item
$f(x)=(x-4)^2+3$
\hint{מצאו את הקודקוד של הפרבולה ובדקו הפיכות בכל צד שלו בנפרד.}
\begin{solbox}
הפונקציה $f(x)=(x-4)^2+3$ אינה חד-חד ערכית  $\mathbb{R}$,
למשל
$f(3)=f(5)=4$.
ניתן לראות שהתמונה שלה שווה
$[3,\infty)$.

היא יורדת ב-$(-\infty,4]$ ועולה ב-$[4,\infty)$. על $[4,\infty)$ ההופכית היא
\[
,f^{-1}(x)=4+\sqrt{x-3}
\]
ועל $(-\infty,4]$ ההופכית היא
\[
.f^{-1}(x)=4-\sqrt{x-3}
\]
\end{solbox}
\item
$f(x)=x^2+2x+2$
\hint{השלימו לריבוע ופעלו כמו בסעיף הקודם.}
\begin{solbox}
נעשה השלמה לריבוע ונקבל כי
הפונקציה
$f(x)=x^2+2x+2=(x+1)^2+1$
אינה חד-חד ערכית
$\mathbb{R}$,
כמו בסעיף הקודם ניתן לראות כי התמונה שלה היא
$[1,\infty)$.

היא יורדת ב-$(-\infty,-1]$
ועולה ב-$[-1,\infty)$. על $[-1,\infty)$ ההופכית היא
\[
,f^{-1}(x)=-1+\sqrt{x-1}
\]
ועל $(-\infty,-1]$ ההופכית היא
\[
.f^{-1}(x)=-1-\sqrt{x-1}
\]
\end{solbox}
\item
$f(x)=\floor{x}$
\begin{solbox}
הפונקציה
$f(x)=\floor{x}$
אינה חד-חד ערכית, ולכן אינה הפיכה  $\mathbb{R}$. היא מונוטונית עולה במובן הרחב, אך לא עולה ממש.
\end{solbox}

\item
$f(x)=x-\floor{x}$
\hint{בכל קטע מהצורה $[n,n+1)$ עם $n$ שלם, כתבו את $f$ בלי סימן הערך השלם.}
\begin{solbox}
הפונקציה $f(x)=x-\floor{x}$ אינה חד-חד ערכית על $\mathbb{R}$, למשל $f(0.5)=f(1.5)$. בכל קטע מהצורה $[n,n+1)$, $n\in\mathbb{Z}$, מתקיים $f(x)=x-n$,
ולכן שם היא עולה ממש וההופכית היא
$f^{-1}(x)=x+n$
על הטווח
$[0,1)$.
\end{solbox}
\item
$f(x)=2^{2^x}$
\hint{כתבו את $f$ כהרכבה של הפונקציה $2^x$ עם עצמה, והשתמשו במונוטוניות שלה.}
\begin{solbox}
נסמן
$g(x)=2^x$
פונקציה זאת עולה ממש והפיכה ולכן ההרכבה
\\
$g(g(x))=2^{(2^x)}$
עולה ממש.
הפונציה ההפוכה היא
$h(x)=log_2(log_2(x))$
.
\end{solbox}
\item
$f(x)=\log_{0.5}(1+x^2)$
\hint{כתבו את $f$ כהרכבה של $1+x^2$ עם $\log_{0.5}$, ושימו לב לכיוון המונוטוניות של כל אחת מהן.}
\begin{solbox}
נסמן
$h(x)=1+x^2$
לא הפיכה על כל הישר.
עולה כאשר
$x>0$
ויורדת כאשר
$x<0$
בכל אחד מהתחומים ניתן להפוך אותה.
\\
כאשר
$x>0$
אז
$x=\sqrt{y-1}$ ,\\
כאשר
$x<0$
אז
$x=-\sqrt{y-1}$.
\end{solbox}
\item
$f(x)=\sinh(x)$
\hint{בדקו את המונוטוניות של כל אחד משני המחוברים בהגדרת $\sinh$. כדי לבודד את $x$, הציבו $t=e^x$ וקבלו משוואה ריבועית ב-$t$.}
\begin{solbox}
$sinh(x)=\frac{e^x-e^{-x}}{2}$
ידוע כי
$\frac{e^x}{2}$
עולה ממש וכי
$\frac{e^{-x}}{2}$
יורדת ממש ולכן
$-\frac{e^{-x}}{2}$
עולה ממש מכאן הפונקציה עולה ממש כסכום של שתי פונקציות עולות ממש.
\end{solbox}
\eMs
\end{enumerate}
